\documentclass{article}
\title{Submersions, Gauss lemma, constant curvature}
\author{Math 6702}
\date{Week 10}

\begin{document}
\maketitle

We connect curvature to geometry in three directions: how curvature behaves under Riemannian submersions, how length and energy behave near a geodesic via Gauss lemma and conjugate points, and how the model spaces of constant curvature drive global theorems like Hadamard-Cartan and Myers.

\section*{Reading}
From \textit{Riemannian Geometry} by Sylvestre Gallot, Dominique Hulin, and Jacques Lafontaine, read \textsection III.D Riemannian Submersions and Curvature: ``Riemannian submersions and connections,'' ``O'Neill's formula,'' and ``Curvature and length of small circles; application to submersions,'' together with \textsection III.E The Behavior of Length and Energy in the Neighborhood of a Geodesic: ``The Gauss lemma'' and ``Conjugate points,'' plus \textsection III.F Manifolds with Constant Sectional Curvature: ``Spheres, Euclidean and hyperbolic spaces,'' and \textsection III.G Topology and Curvature: ``The Myers and Hadamard-Cartan theorems.''

\end{document}
